% ════════════════════════════════════════════
% 诊断指标与判别标准（5.1.3）
% ════════════════════════════════════════════
\subsubsection{诊断指标与判别标准}

本节从数据结构、分布形态与变异来源三个角度建立数据理解框架，把定性判断转化为可复核的数值证据。

在数据结构上，把 500 条帖子视为同一主页的 500 次独立观测，每条记录包含 7 个输入变量与 1 个曝光结果。计量单位上，X1 以人为单位、Y1 以次为单位，X4、X5、X6 分别以月、星期与小时计，X2、X3、X7 为类别编码。按帖子发布前运营人员能否决定，变量可分为两类：帖子类型、内容主题、月份、星期、时段与付费标记属于发帖决策变量；粉丝规模由历史积累决定，属于账号状态变量；曝光总次数是发布后的结果。数据没有帖子编号、发布时间戳与重复观测，一行即一条帖子，因而判定为帖子级截面数据，而不是面板或时间序列。月份与星期只刻画发布的周期性位置，不构成对同一条帖子的重复跟踪。

在分布形态上，曝光量为非负计数，需先判断尾部厚度。本节计算样本偏度与超额峰度，前者衡量分布围绕均值的非对称程度，后者衡量尾部相对于正态分布的厚度；偏度明显大于 1 时分布右偏，超额峰度明显大于 0 时尾部偏厚。同时计算最大值与中位数之比，作为极端值相对规模的直观度量，并比较原始曝光量与 log1p 曝光量的偏度、峰度，判断变换是否有效压缩尾部。

在关联结构上，对粉丝规模与月份分别计算 Pearson 线性相关与 Spearman 秩相关。前者度量二者线性同步的程度，后者度量单调同步的程度；两个系数同时接近 1，说明粉丝规模几乎随月份单调累积，而不是两个可以独立变动的变量。

在变异来源上，进一步建立粉丝规模对数对月份哑变量的辅助回归。该回归的解释度衡量粉丝规模变异中有多少可由月份解释，月内方差占比则是 1 减去该解释度；解释度越接近 1、月内占比越接近 0，说明粉丝规模的变异几乎全部发生在月间，月内只剩极少量信息。

综合上述指标，本节采用一组明确的判别标准：若 Pearson 相关系数绝对值不低于 0.9、Spearman 相关系数不低于 0.9、月份对粉丝规模对数的解释度不低于 0.9，且月内方差占比低于 5\%，则判定粉丝规模与月份属于同一变异维度，同时进入回归会造成系数不可识别。该标准把“看起来相关”提升为“变异来源重合”的判断，避免仅凭相关系数下结论。

综上所述，本节建立了由数据结构判定、尾部度量、关联度量与方差分解组成的数据理解框架，得到了判断变量可识别性的量化标准，接下来将基于此对 500 条帖子进行求解。
